%% thermalops_manuscript.tex
%% SIGMETRICS/POMACS 2027 submission — double-anonymous review
%% ACM single-column format
\documentclass[acmsmall,screen,review,anonymous]{acmart}

\usepackage{booktabs}
\usepackage{multirow}
\usepackage{microtype}
\usepackage{xspace}
\usepackage{amsmath}
\usepackage{siunitx}

%% Suppress ACM copyright/ISBN blocks for review
\setcopyright{none}
\acmDOI{}
\acmISBN{}
\acmConference[SIGMETRICS '27]{ACM SIGMETRICS}{2027}{ACM}
\acmYear{2027}

%% ----------------------------------------------------------------
%% Shorthand macros
%% ----------------------------------------------------------------
\newcommand{\Frontier}{\textsc{Frontier}\xspace}
\newcommand{\dlo}{\ensuremath{\Delta_{\mathrm{lo}}}\xspace}
\newcommand{\dhi}{\ensuremath{\Delta_{\mathrm{hi}}}\xspace}
\newcommand{\Bref}{\ensuremath{B_{\mathrm{ref}}}\xspace}
\newcommand{\Sref}{\ensuremath{S_{\mathrm{ref}}}\xspace}
\newcommand{\Cl}{\ensuremath{C_L}\xspace}
\newcommand{\Cu}{\ensuremath{C_U}\xspace}
\newcommand{\pp}{\,pp\xspace}
\newcommand{\eg}{\textit{e.g.},\xspace}
\newcommand{\ie}{\textit{i.e.},\xspace}
\newcommand{\cf}{\textit{cf.}\xspace}

%% ----------------------------------------------------------------
\begin{document}

\title{Persistent Node-Level Thermal Fingerprints on Frontier:\\
       Measurement and Constraint-Aware Replay}

\begin{abstract}
GPU supercomputers exhibit persistent node-level thermal heterogeneity: nodes that run
warm during one workload tend to run warm in subsequent workloads, independent of job
composition.  We characterize this property---which we call a \emph{thermal
fingerprint}---across 2,637 production workloads on the 9,856-node \Frontier system.
Fingerprints learned from 2024 training telemetry prospectively predict each node's mean
relative thermal position in 2025 workloads with Pearson $r = 0.926$, Spearman
$\rho = 0.936$, and direct no-refit $R^2 = 0.848$ at the node level.

We then quantify the scheduling-capacity advantage of fingerprint-aware placement through
a preregistered held-out constraint-aware replay study.  All comparison hypotheses,
estimators, and analysis rules are frozen at a fixed commit before any scenario executes.
A 25-scenario held-out set evaluates eight comparisons under a factorial placebo structure.
Under a concurrent-load cabinet reference at the preregistered 1.25$\times$ choice ratio,
fingerprint-guided placement yields a median conservative bracket separation of
$+12.48\pp$ versus random placement (bootstrap 95\% CI $[+11.09,\;+19.04]\pp$; 25/25
scenarios; Holm-adjusted $p = 4.77 \times 10^{-7}$).  The advantage survives three
independent controls: a within-cabinet permutation placebo ($+10.88\pp$; 24/25), a global
permutation placebo ($+12.31\pp$; 25/25), and a leave-one-out structural-position template
($+13.95\pp$; 25/25), collectively indicating that individual node identity contributes
beyond cabinet membership and fixed rack geometry.  Random thermal spreading provides no
statistically distinguishable benefit over random placement.  The direction of the
fingerprint advantage is preserved under thermal threshold definitions at \SI{61.875}{\celsius},
\SI{63.625}{\celsius}, and \SI{66.6625}{\celsius}.  Results are bounded by the replay
model's scope and cannot be extrapolated to realized production throughput gains.
\end{abstract}

%%
\keywords{thermal fingerprinting, GPU supercomputers, scheduling, constraint-aware replay,
          HPC measurement, reproducibility}

\maketitle

%% ================================================================
\section{Introduction}
\label{sec:intro}
%% ================================================================

Modern GPU supercomputers are thermally heterogeneous at the node level.  Two nodes
allocated to the same job at the same time can exhibit systematically different thermal
behaviors even when running identical workloads, and this heterogeneity is not random
noise: it is \emph{persistent}.  A node that ran warm in 2024 tends to run warm in 2025.
This property---which we term a \emph{thermal fingerprint}---is stable enough to
characterize from historical telemetry and potentially exploit in placement decisions.

Despite its operational relevance, node-level thermal fingerprinting at production scale
remains undercharacterized.  Prior thermal-aware scheduling work~\cite{fan2007power,
inadomi2015hpc,ostrouchov2020machine,sinha2023thermal} has addressed temperature-aware
placement or workload diversification, but has not measured or exploited \emph{per-node
persistent fingerprints} learned prospectively from historical telemetry at the scale of
a leadership-class supercomputer.

This paper asks two questions.  First, how stable and predictive are node-level thermal
fingerprints on a leadership-class GPU supercomputer?  Second, if such fingerprints exist,
can a thermal-fingerprint-aware placement policy exploit them to increase scheduling
capacity yield under realistic thermal operating constraints?

To answer these questions we use the \Frontier system~\cite{frontier_userguide}---a
9,856-node, 9.2 exaflop GPU supercomputer---and its publicly available
telemetry~\cite{frontier_olcf_dataset}.  We measure node-level thermal fingerprints from
3,945 training jobs spanning 2024 and evaluate their prospective predictive accuracy on
2,637 test jobs in 2025.  We then design and execute a preregistered held-out replay
study that evaluates eight scheduling policies under a constraint-aware simulation that
replays the full base workload before measuring incremental placement capacity.

\paragraph{Contributions.}
\begin{enumerate}
  \item \textbf{Large-scale fingerprint characterization.}  We characterize persistent
    node-level thermal heterogeneity across a 2,000-node cohort spanning 77 cabinets of
    \Frontier over two calendar years.  Fingerprints computed from 2024 telemetry
    prospectively predict 2025 thermal positions at the node level with Pearson $r = 0.926$,
    Spearman $\rho = 0.936$, and direct no-refit $R^2 = 0.848$---without any intercept or
    slope adjustment.

  \item \textbf{Validated constraint-aware replay model.}  A stateful replay model
    reconstructs cabinet-level thermal budgets from raw telemetry and replays the full
    historical base workload before measuring incremental placement capacity yield under
    thermal operating constraints.  Cabinet power estimation MAE is \SI{0.37}{\kilo\watt}
    (0.8\% of the P99 spread across cabinets).

  \item \textbf{Preregistered factorial challenge design.}  All hypotheses, the estimator,
    bootstrap parameters, and interpretation rules are frozen at a fixed commit before
    execution.  A 16-condition preflight gate---including simulation equivalence at
    $10^{-9}$ tolerance, regression tests, and seed-index verification---must pass before
    any scenario runs.

  \item \textbf{Node-specificity placebo structure.}  Eight preregistered comparisons
    isolate the contribution of individual node identity from cabinet membership
    (within-cabinet permutation placebo), arbitrary global ranking (global permutation
    placebo), spreading mechanics (FP-guided spreading vs.\ random spreading), and a
    leave-one-out structural position template.

  \item \textbf{Thermal threshold sensitivity.}  A post-exposure sensitivity analysis
    evaluates the FP-vs-Random direction under three threshold definitions (P90, P95,
    P99).  The advantage sign is preserved at all three, providing directional robustness
    evidence.
\end{enumerate}

\paragraph{Paper organization.}
Section~\ref{sec:background} surveys related work.
Section~\ref{sec:fingerprint} characterizes the thermal fingerprint.
Section~\ref{sec:model} describes the replay model.
Section~\ref{sec:protocol} describes the preregistered challenge protocol.
Section~\ref{sec:results} presents the primary results.
Section~\ref{sec:robustness} presents robustness and sensitivity analyses.
Section~\ref{sec:discussion} discusses limitations and scope.
Section~\ref{sec:conclusion} concludes.

%% ================================================================
\section{Background and Related Work}
\label{sec:background}
%% ================================================================

\paragraph{Thermal management in HPC.}
Fan, Weber, and Barroso~\cite{fan2007power} documented the importance of power and
cooling efficiency in warehouse-scale computing.  Thermal-aware job scheduling has been
studied in the context of rack-level temperature constraints~\cite{inadomi2015hpc},
workload diversification to reduce hotspot formation~\cite{ostrouchov2020machine}, and
opportunistic admission based on real-time thermal headroom~\cite{sinha2023thermal}.
These approaches rely on \emph{reactive} measurements of current thermal state rather than
\emph{prospective} node-level fingerprints learned from historical telemetry.

\paragraph{Node-level heterogeneity.}
Hardware heterogeneity within nominally identical node classes is well documented in
HPC~\cite{ostrouchov2020machine}.  At the node level, persistent thermal differences can
arise from manufacturing variation, placement within the cooling circuit, BIOS
configuration, and local airflow patterns.  Because system hostnames and physical
locations are fixed (\ie a 1:1 mapping exists between hostname and hardware slot),
persistent thermal fingerprints may reflect node-specific hardware characteristics,
fixed slot-level cooling properties, or both; the two cannot be disentangled from
telemetry alone (Section~\ref{sec:discussion}).

\paragraph{Thermal-aware scheduling systems.}
TAPAS~\cite{tapas} and PAL~\cite{pal} integrate thermal feedback into admission or
placement decisions.  These systems operate on real-time thermal signals; in contrast,
our approach uses a fingerprint computed offline from historical telemetry, which is
available at scheduling time without instrumentation overhead.

\paragraph{Preregistered experimental design in systems research.}
Preregistration---specifying hypotheses, analysis rules, and estimators before collecting
outcome data---is well established in behavioral and clinical research.  Its adoption in
systems measurement is more recent~\cite{mytkowicz2009producing}.  We implement a strict
variant: the scientific runner, all comparison definitions, the estimator, bootstrap seed,
and interpretation rules are all committed to a version-controlled repository at a fixed
commit before any scenario is executed.  This design separates specification from
execution and provides verifiable, machine-checkable provenance for all results.

\paragraph{Existing Frontier measurement literature.}
ExaDigiT~\cite{exadigit} provides a digital-twin simulation capability for Frontier.
Our replay model is independent: it uses the publicly available job telemetry dataset
from OLCF~\cite{frontier_olcf_dataset} and does not rely on the ExaDigiT infrastructure.

%% ================================================================
\section{Thermal Fingerprint Characterization}
\label{sec:fingerprint}
%% ================================================================

\subsection{System and Telemetry}

\Frontier is a 9,856-node AMD-GPU-based supercomputer at Oak Ridge National Laboratory
with a peak performance of 9.2 exaflops (FP64).  Nodes are organized in 77 cabinets
with dedicated liquid cooling per cabinet.  The per-node thermal metric we use is the
job-relative departure from the node's expected thermal behavior: for job $j$ and
node $n$, the fingerprint target is the \emph{median} job-relative
$\Delta T_{\mathrm{p95}}$, defined as the node's p95 GPU temperature minus the job-wide
p95 GPU temperature, aggregated across all measurements within the job.

We use the publicly available \Frontier job metadata and per-node GPU temperature
telemetry~\cite{frontier_olcf_dataset}.  The training set consists of 3,945 jobs
completed in 2024; the held-out test set consists of 2,637 jobs completed in 2025.
No information from the 2025 test set is used to compute fingerprint scores.

\subsection{Fingerprint Definition and Temporal Stability}

The \emph{thermal fingerprint} of node $n$ is the median job-relative $\Delta T_{\mathrm{p95}}$
computed over all 2024 jobs for which node $n$ participated.  We evaluate two
2024 cohorts---\textbf{FA} ($n=1,976$ nodes) and \textbf{FC} ($n=1,466$ nodes)---and
the combined prospective test.

At the node level, across 64,545 node-job observations in 2025, fingerprints computed
from 2024 telemetry predict 2025 job-relative $\Delta T_{\mathrm{p95}}$ with:
\begin{itemize}
  \item Node-level Pearson $r = 0.926$ ($p < 10^{-300}$)
  \item Node-level Spearman $\rho = 0.936$
  \item Direct no-refit $R^2 = 0.848$ (no slope or intercept adjustment applied)
\end{itemize}
Node-level sign agreement (fingerprint and 2025 deviation share the same sign) is
$90.7\%$.  The median absolute prediction error at the node level is \SI{0.26}{\celsius}.

\paragraph{Why node-level metrics.}
Job-level correlations ($r \approx 0.45$, $R^2 \approx 0.21$) are lower because a
single job may span nodes with varying fingerprints, diluting the signal.  The
node-level summary---using each node's mean 2025 position across all its test
jobs---is the appropriate target for a placement policy that selects specific nodes.

\paragraph{Robustness.}
The predictive relationship holds across power-quartile strata, across both 2024 cohorts
(FA and FC), and for both p95 temperature and mean temperature as alternative targets
(Supplemental material).  Null-permutation testing ($n=1,000$ permutations) confirms
the observed ICC for both cohorts exceeds the 99th percentile of the null distribution
($p_{\mathrm{perm}} < 0.001$).

% [FIGURE 1: Node-level fingerprint distribution and temporal stability scatter plot.
%  Two panels: (A) histogram of 2024 fingerprint scores (FA + FC cohorts);
%  (B) scatter of 2024 fingerprint score vs. 2025 mean node-relative dT_p95,
%  with r, rho, R^2 annotated.  Source: phase3a_fingerprint_2024.parquet,
%  phase3a_test_set.parquet.]

\subsection{Causal Caveat}

Because \Frontier assigns a 1:1 mapping between hostname and physical chassis slot,
persistent thermal heterogeneity could reflect node-specific hardware characteristics
(\eg GPU manufacturing variation, thermal paste application) or fixed slot-level
cooling properties (\eg position in the hydraulic cooling circuit), or both.  Telemetry
alone cannot distinguish these; we use the fingerprint as a predictive score without
claiming to identify its physical origin.

%% ================================================================
\section{Constraint-Aware Replay Model}
\label{sec:model}
%% ================================================================

\subsection{Model Overview}

The replay model measures the incremental scheduling capacity yield of a placement
policy under realistic thermal operating constraints.  It proceeds in two phases:

\begin{enumerate}
  \item \textbf{Base workload replay.}  The full historical base workload (all background
    jobs in scenario $s$) is replayed in original order and timing, producing a reference
    thermal state for each cabinet.

  \item \textbf{Incremental probe.}  A synthetic incremental workload is then placed
    at varying load factors $\lambda$.  For each $\lambda$, the model evaluates whether
    the cabinet's thermal operating envelope is exceeded.  The fraction of incremental
    load that can be absorbed before envelope exceedance defines the capacity yield $C(\lambda)$.
    A binary search over $\lambda$ identifies the crossing interval $[C_L, C_U]$.
\end{enumerate}

\subsection{Cabinet Thermal Reference}

We use a \emph{concurrent-load cabinet reference}: the reference thermal state for a
cabinet is derived from the concurrently scheduled jobs that actually ran alongside
the evaluated workload, not from the evaluated workload itself.  This avoids a
circularity in the original job-contribution reference (\cf Section~\ref{sec:hier}),
where the job being evaluated partially defined the reference it was measured against.

The thermal exposure metrics are:
\begin{align}
  B_{\mathrm{ref},s} &= \text{fraction of time at which cabinet thermal load exceeds threshold } T \\
  S_{\mathrm{ref},s} &= \text{mean thermal load during exceeding intervals}
\end{align}
Both are computed from concurrently scheduled jobs only, holding fixed the threshold $T$
(Section~\ref{sec:robustness} evaluates three threshold definitions).

\subsection{Model Validation}

Cabinet power estimation MAE is \SI{0.37}{\kilo\watt}, representing 0.8\% of the P99
spread of cabinet-level power across the deployment.  Spearman correlation between
modeled and actual cabinet thermal burden exceeds 0.97 across validation scenarios.

% [FIGURE 2: Capacity yield crossing illustration.
%  Panel (A): lambda sweep showing C(lambda) for one scenario/policy pair,
%  with C_L and C_U annotated.  Panel (B): delta_lo = C_L(FP) - C_U(Random)
%  illustrated for one scenario, showing verdict A.]

\subsection{Conservative Interval Comparison Estimand}

For policies $A$ and $B$ and scenario $s$, define:
\begin{equation}
  \dlo(A,B) = C_L(A) - C_U(B), \quad
  \dhi(A,B) = C_U(A) - C_L(B).
\end{equation}
The scenario verdict is:
\begin{itemize}
  \item \textbf{A} if $\dlo > 0$ ($A$'s lower bound exceeds $B$'s upper bound);
  \item \textbf{B} if $\dhi < 0$ ($B$'s lower bound exceeds $A$'s upper bound);
  \item \textbf{I} (indeterminate) otherwise.
\end{itemize}
The conservative bracket ensures that \textbf{A} verdicts reflect a genuine non-overlapping
advantage: no assignment of values within each interval would yield a different conclusion.

%% ================================================================
\section{Preregistered Challenge Protocol}
\label{sec:protocol}
%% ================================================================

\subsection{Preregistration and Freeze}

The scientific runner (SHA256-verified), all comparison definitions, the estimator,
bootstrap seed and draw count, and interpretation rules are committed to a version-controlled
repository at a fixed commit before any scenario is executed.  The 25 held-out scenarios
(sc30--sc54) are not touched until the preregistration commit is recorded.

\subsection{Preflight Gate}

Execution requires passing 16 preflight conditions, including:
simulation equivalence at $10^{-9}$ tolerance; bootstrap regression test (seed 9999,
5000 draws, expected CI within tolerance); seed-index correctness for all scenarios;
SHA256 verification of reference files; and clean git tracked state.  The gate must
pass in full before scenario dispatch begins.

\subsection{Policies and Placebo Structure}

Seven policies are evaluated at choice ratio $r = 1.25$:

\begin{description}
  \item[\textbf{FP}] Full fingerprint selection: among $\lfloor r \cdot k \rfloor$
    candidate nodes, select those with the lowest (coolest) fingerprint scores.
  \item[\textbf{Random}] Uniform random selection from candidates (no fingerprint
    information).
  \item[\textbf{Location}] Leave-one-out structural position comparator.
    For each job in the held-out set, a score is constructed from all \emph{other}
    jobs sharing the same physical position, excluding the evaluated job itself.
    The fallback hierarchy applies level-1 (exact-position LOO, 1,990/2,000 nodes),
    level-3 (adjacent-level LOO, 3 nodes), and global fallback (7 nodes) as needed;
    no node receives a NaN score.  The resulting template correlates with the
    ground-truth structural template at Spearman $\rho = 0.9823$.
  \item[\textbf{FP-Spread}] Fingerprint-guided spreading: spread candidates across
    fingerprint-score quantiles to maintain thermal diversity while still using
    fingerprint information.
  \item[\textbf{Rand-Spread}] Random spreading: spread candidates by random
    assignment; no fingerprint information.
  \item[\textbf{FP-Global-Placebo}] Global fingerprint permutation: fingerprint
    scores randomly permuted across the entire node population before selection.
  \item[\textbf{FP-WithinCab-Placebo}] Within-cabinet fingerprint permutation:
    scores permuted within each cabinet, preserving cabinet-level structure.
\end{description}

\paragraph{Placebo invariant.}
In all placebo conditions, the thermal reconstruction after placement always uses
each selected node's \emph{true frozen fingerprint score}, not the permuted score.
The placebo score affects only which nodes are selected; thermal accounting uses
ground-truth node properties.

\subsection{Bootstrap and Sign Test}

For each comparison, we compute:
\begin{itemize}
  \item Median $\dlo(A,B)$ and $\dhi(A,B)$ across the 25 scenarios;
  \item Bootstrap 95\% CI (5,000 paired scenario draws, seed 9999, percentile method);
  \item Sign test $p$-value on decisive (A or B) scenarios using a two-sided exact
    binomial test under the null hypothesis of equal probability;
  \item Holm--Bonferroni adjustment across all eight comparisons, $\alpha = 0.05$.
\end{itemize}

\paragraph{Bootstrap scope.}
The 25 replay scenarios draw from the same underlying 2,637-job 2025 base workload.
Bootstrap intervals therefore characterize variation across the defined replay scenarios
conditional on that historical workload, not across independently sampled years or
production deployments.

\subsection{Experimental Hierarchy}
\label{sec:hier}

\paragraph{Original analyses.}
An initial choice-ratio experiment (Experiment~1) and a node-specificity challenge
(Experiment~2) were pre-specified sequentially under a \emph{job-contribution cabinet
reference}, in which the evaluated job's own thermal behavior contributed to the cabinet
reference against which it was measured.

\paragraph{Reference correction.}
A review of the reference construction identified a circularity: the evaluated job's
scheduling period partially overlapped with the collection window for the reference.
We corrected this by adopting the \emph{concurrent-load cabinet reference}
(Section~\ref{sec:model}), which uses only jobs concurrently scheduled alongside the
evaluated workload.

\paragraph{Post-exposure analyses.}
All corrected-reference analyses reported in this paper---the choice-ratio results
(Section~\ref{sec:results}), node-specificity challenge (Section~\ref{sec:results}),
threshold sensitivity (Section~\ref{sec:robustness}), and late-cabinet diagnostic
(Section~\ref{sec:robustness})---were performed after the original results were already
observed.  These analyses are post-exposure reference-sensitivity and robustness analyses;
they are not retroactively preregistered primary experiments.  The original results are
reported in Section~\ref{sec:original} for completeness.

%% ================================================================
\section{Primary Results}
\label{sec:results}
%% ================================================================

\subsection{Original Results (Job-Contribution Reference)}
\label{sec:original}

Under the original job-contribution cabinet reference at $r = 1.25$, FP placement
yielded a median $\dlo$ of $+6.71\pp$ versus random placement (22/25 scenarios;
bootstrap 95\% CI $[+3.92,\;+8.72]\pp$).  The within-cabinet placebo separation
was $+3.54\pp$ (20/25; CI $[+0.56,\;+4.44]\pp$).  FP-guided spreading outperformed
direct FP at the highest choice ratios ($r \geq 8$, unrestricted), consistent with
a constraint-regime where spreading is capacity-increasing when the cabinet envelope
is not yet binding on individual node selection.

These values are not the primary claims; they are reported for completeness and to
document the effect of the reference correction.

\subsection{Corrected-Reference Choice-Ratio Experiment (Experiment~1)}
\label{sec:exp1}

Under the concurrent-load cabinet reference, we evaluate all seven choice ratios:
$r \in \{1.0, 1.25, 1.5, 2, 4, 8, \text{unrestricted}\}$.

% [TABLE 1: Experiment 1 — verdict summaries (n_A/n_I/n_B) and median delta_lo
%  for FP vs Random, FP vs Location, FP vs FP-Spread, FP-Spread vs Random
%  across all seven choice ratios.  Source: c1_correctedref_analysis.json.]

Key findings under the corrected reference:
\begin{itemize}
  \item FP placement beats Random at every ratio with actual choice ($r \geq 1.25$).
    The $r = 1.0$ condition is a negative control (no candidate diversity) and
    yields an indeterminate result as expected.
  \item The Location comparator (structural position LOO) is a strong baseline and
    must be held visible alongside FP; FP's numerical advantage over Location is more
    modest than its advantage over Random.
  \item Direct FP outperforms FP-guided spreading under the corrected reference
    ($\dlo(\text{FP-Spread},\text{FP}) = -9.18\pp$; 0/0/25; Holm $p = 4.77\times10^{-7}$).
    This reverses the high-choice-ratio FP-Spread advantage observed under the original
    reference, because the concurrent-load reference changes the binding constraint
    regime: spreading becomes capacity-reducing when the active constraint is thermal
    incidence of the evaluated job rather than cabinet envelope headroom.
\end{itemize}

\paragraph{Interpreting the reference correction.}
The change from the job-contribution to the concurrent-load reference does not simply
``relax'' or ``tighten'' one constraint uniformly.  It changes \emph{which} constraint
regime is active.  Under the job-contribution reference, the evaluated job's own thermal
contribution inflated the apparent cabinet reference, making spreading appear beneficial
by distributing that contribution across cabinets.  Under the concurrent-load reference,
the cabinet load is externally determined; the dominant constraint becomes the ratio of
the fingerprint-matched node's thermal exposure to the concurrent reference.  Numerical
deltas therefore cannot be compared directly across the two reference definitions as
evidence of a ``stronger'' or ``weaker'' effect; they reflect different physical regimes.

\subsection{Corrected-Reference Node-Specificity Challenge (Experiment~2)}
\label{sec:exp2}

Table~\ref{tab:exp2} reports all eight comparisons at $r = 1.25$ under the corrected
reference.  This is the most mechanistically informative experiment: the factorial
placebo structure tests whether the fingerprint advantage can be attributed to
cabinet membership, global ranking, spreading mechanics, or structural position,
or whether individual node identity is required.

\begin{table}[t]
\caption{Node-specificity challenge results at $r = 1.25$ under the concurrent-load
cabinet reference.  $\dlo$ and CI in percentage points.
$p_\mathrm{Holm}$: Holm-adjusted sign-test $p$-value; 8 comparisons, $\alpha = 0.05$.}
\label{tab:exp2}
\begin{tabular}{lrrlrr}
\toprule
Comparison & $n_A/n_I/n_B$ & Med.\ $\dlo$ & 95\% CI & $p_\mathrm{Holm}$ & Sig. \\
\midrule
FP vs.\ Random       & 25/0/0  & $+12.48$ & $[+11.09,\;+19.04]$ & $4.77\!\times\!10^{-7}$ & Yes \\
FP vs.\ WithinCab   & 24/1/0  & $+10.88$ & $[+5.96,\;+17.40]$  & $4.77\!\times\!10^{-7}$ & Yes \\
FP vs.\ GlobalPlacebo & 25/0/0 & $+12.31$ & $[+10.88,\;+17.40]$ & $4.77\!\times\!10^{-7}$ & Yes \\
FP vs.\ Location     & 25/0/0  & $+13.95$ & $[+12.23,\;+21.98]$ & $4.77\!\times\!10^{-7}$ & Yes \\
FP-Spread vs.\ Rand-Spread & 23/2/0 & $+4.42$ & $[+2.14,\;+5.42]$ & $4.77\!\times\!10^{-7}$ & Yes \\
FP-Spread vs.\ FP   & 0/0/25  & $-9.18$  & $[-17.91,\;-6.60]$  & $4.77\!\times\!10^{-7}$ & Yes \\
Rand-Spread vs.\ Random & 5/14/6 & $-0.19$ & $[-0.58,\;0.00]$  & $1.000$                 & No  \\
FP-Spread vs.\ Location & 24/1/0 & $+4.96$ & $[+2.26,\;+8.34]$  & $4.77\!\times\!10^{-7}$ & Yes \\
\bottomrule
\multicolumn{6}{l}{\footnotesize Bootstrap: 5,000 paired draws, seed 9999, percentile 95\% CI.} \\
\multicolumn{6}{l}{\footnotesize Holm--Bonferroni: 8 comparisons, $\alpha = 0.05$.}
\end{tabular}
\end{table}

% [FIGURE 3: Forest plot of 8 comparisons at r=1.25, corrected reference.
%  x-axis: median delta_lo (pp), whiskers: 95% CI.
%  Color: green=significant, grey=nonsignificant.
%  Source: c2_correctedref_analysis.json.]

\paragraph{Mechanistic interpretation.}
The pattern in Table~\ref{tab:exp2} allows the following qualified statements, bounded
to the replay model's scope:

\begin{itemize}
  \item \textbf{Cabinet membership does not explain the advantage} (FP vs.\ WithinCab:
    24/1/0, $+10.88\pp$).  Permuting fingerprint scores within each cabinet---preserving
    the cabinet-level score distribution---materially degrades performance, indicating
    that individual node identity within the cabinet contributes beyond cabinet-level
    composition.

  \item \textbf{Generic global ranking does not explain the advantage} (FP vs.\
    GlobalPlacebo: 25/0/0, $+12.31\pp$).  A globally permuted placebo---which retains
    the same score distribution but destroys within-population rank order---provides no
    benefit, ruling out ranking artifacts.

  \item \textbf{Fixed rack position does not fully explain the advantage} (FP vs.\
    Location: 25/0/0, $+13.95\pp$).  The structural leave-one-out position template
    materially underperforms the thermal fingerprint, consistent with the fingerprint
    capturing information beyond what historical position assignment alone explains.
    However, because hostnames and slots are 1:1 mapped, this comparison cannot rule
    out node-specific hardware heterogeneity as a contributor (Section~\ref{sec:discussion}).

  \item \textbf{Random spreading provides no benefit} (Rand-Spread vs.\ Random:
    5/14/6, Holm $p = 1.0$).  Thermal spreading without fingerprint information
    yields an indeterminate result consistent with null effect.

  \item \textbf{Fingerprint ranking helps even under spreading} (FP-Spread vs.\
    Rand-Spread: 23/2/0, $+4.42\pp$).  Fingerprint information improves outcome even
    when the spreading constraint is applied.

  \item \textbf{Direct FP outperforms FP-guided spreading} (FP-Spread vs.\ FP:
    0/0/25, $-9.18\pp$).  Concentration, not dispersion, is the optimal strategy under
    the concurrent-load reference regime.
\end{itemize}

\paragraph{Left-boundary behavior.}
Random placement frequently encounters left-boundary scenarios (no crossing found within
the lambda sweep) under the corrected reference at $r = 1.25$.  This reflects that many
scenarios have essentially zero incremental capacity under random placement when the
concurrent cabinet load is already near the thermal operating envelope.  FP encounters
zero left-boundary scenarios at $r = 1.25$.  This floor behavior makes the absolute
median $\dlo$ values conservative estimates of the true separation: in left-boundary
scenarios, Random's capacity is at most zero, while FP's is positive.

% [FIGURE 4: Absolute capacity distributions for FP and Random at r=1.25.
%  Violin or box plots of C_L per policy, with left-boundary count annotated.
%  Source: c2_correctedref_analysis.json.]

%% ================================================================
\section{Robustness and Sensitivity}
\label{sec:robustness}
%% ================================================================

\subsection{Thermal Threshold Sensitivity}
\label{sec:thresh}

The threshold $T$ used to define thermal incidence ($B_{\mathrm{ref}}$, $S_{\mathrm{ref}}$)
is a parameter of the analysis.  We evaluate three threshold definitions drawn from
the empirical distribution of node-level peak GPU temperatures on \Frontier:
P90 = \SI{61.875}{\celsius},
P95 = \SI{63.625}{\celsius} (the preregistered threshold), and
P99 = \SI{66.6625}{\celsius}.

\paragraph{Important: these are alternative exposure definitions, not nested constraints.}
At each threshold, both $B_{\mathrm{ref},s}$ and $S_{\mathrm{ref},s}$ are recomputed
independently.  A higher threshold does not simply ``relax'' the P95 constraint;
it defines a different historical thermal exposure against which policies are evaluated.
Do not interpret the magnitude progression as a monotonic dose-response to a loosening
constraint.

\begin{table}[t]
\caption{Threshold sensitivity: FP vs.\ Random under three threshold definitions.
$\dlo$ and CI in percentage points.  Random left-boundary counts (lb) shown.}
\label{tab:thresh}
\begin{tabular}{lrlrr}
\toprule
Threshold & $n_A/n_I/n_B$ & Med.\ $\dlo$ & 95\% CI & Random lb \\
\midrule
P90 (\SI{61.875}{\celsius}) & 24/1/0 & $+8.90$ & $[+4.95,\;+14.50]$ & 3/25 \\
P95 (\SI{63.625}{\celsius}) & 25/0/0 & $+12.48$ & $[+11.09,\;+19.04]$ & 9/25 \\
P99 (\SI{66.6625}{\celsius}) & 25/0/0 & $+21.33$ & $[+16.39,\;+25.36]$ & 19/25 \\
\bottomrule
\multicolumn{5}{l}{\footnotesize P95 result reuses the corrected-reference analysis from Section~\ref{sec:results}.} \\
\multicolumn{5}{l}{\footnotesize P90 and P99 are post-exposure sensitivity probes. Bootstrap and seed identical to Table~\ref{tab:exp2}.}
\end{tabular}
\end{table}

\paragraph{Primary conclusion.}
The sign of the FP advantage is preserved under all three threshold definitions;
the direction does not reverse at either endpoint (Table~\ref{tab:thresh}).
The P99 magnitude ($+21.33\pp$) is substantially floor-limited: at P99, 19 of 25
Random scenarios are left-boundary limited, meaning Random has near-zero incremental
capacity.  The large P99 estimate reflects Random's floor behavior at the hottest
threshold, not a larger true FP advantage.

\paragraph{Binding constraint shift.}
The first binding constraint for FP placement shifts between threshold definitions:
at P90 and P95, all FP crossings are driven by the thermal-incidence ratio ($B_\mathrm{ratio}$:
the ratio of actual incidence to $B_{\mathrm{ref},s}$ reaches the capacity boundary);
at P99, 20 of 25 FP crossings are driven by the thermal-severity ratio ($S_\mathrm{ratio}$:
the ratio of actual severity to $S_{\mathrm{ref},s}$ reaches the boundary first).
Both are thermal constraints; the shift reflects the different exposure characterization
at the hotter threshold.

\subsection{Late-Cabinet Diagnostic}

Three late-observed cabinets (installed partway through the data collection period)
could in principle dominate the analysis if their thermal behavior differs systematically
from earlier cabinets.  We audit node-hours, event-touches, and FP-vs-Random enrichment
for these three cabinets specifically.  The diagnostic does not reveal that these cabinets
dominate the result; their structural share of the held-out workload is proportional to
their size in the active fleet.

A limitation of the concurrent-load reference: jobs crossing collection-day boundaries
may be incompletely observed in the concurrent scheduling window, so some apparent
zero-load periods may reflect unavailable telemetry rather than true zero load.  We do
not quantify a specific bias magnitude from this effect, but flag it as a potential
source of underestimation for the reference load.

%% ================================================================
\section{Discussion and Limitations}
\label{sec:discussion}
%% ================================================================

\paragraph{Replay-model scope.}
All results are obtained from a constraint-aware replay model, not from live deployment
on \Frontier.  The replay model measures incremental placement capacity yield under
thermal constraints given the historically observed workload; it cannot be extrapolated
to realized production throughput gains, which depend on queue dynamics, job arrival
rates, user behavior, and system management policies not captured in the replay.

\paragraph{Hardware-versus-fixed-slot causality.}
Because hostnames and physical chassis slots are 1:1 mapped on \Frontier, persistent
thermal heterogeneity could reflect node-specific hardware properties (\eg GPU die
characteristics, thermal interface materials) or fixed slot-level cooling properties
(\eg position in the hydraulic manifold).  The within-cabinet and position-LOO controls
provide evidence that individual node identity contributes beyond cabinet membership and
beyond the self-excluding structural position template, but they cannot resolve
hardware-versus-slot causality.  An experiment with deliberate node migrations between
slots would be required to separate these effects, which is outside the scope of this
study.

\paragraph{Bootstrap scope and independence.}
The 25 replay scenarios draw from the same underlying 2,637-job historical base workload.
Bootstrap intervals characterize variation across the defined replay scenarios conditional
on that workload; they are not confidence intervals from independently sampled
production systems or independently sampled years.

\paragraph{Position-LOO construction.}
The structural position template is constructed by a leave-one-out procedure excluding the
evaluated job from its own position's history.  The fallback hierarchy applies when a
node lacks sufficient self-excluding history: 1,990 of 2,000 nodes receive a level-1
exact-position LOO score; 3 receive level-3 scores; 7 receive a global fallback score;
no node receives a NaN score.  The seven global fallbacks were retained under an explicit
pre-run protocol decision.  The template Spearman correlation with the ground-truth
structural template is $\rho = 0.9823$.

\paragraph{Experimental hierarchy and post-exposure analyses.}
The corrected-reference analyses in Sections~\ref{sec:results} and~\ref{sec:robustness}
were performed after the original results were observed.  The reference correction is
methodologically motivated---the job-contribution reference has a structural circularity
---but the corrected-reference analyses are not retroactively preregistered primary
experiments.  We report the original results (Section~\ref{sec:original}) alongside
the corrected-reference analyses and do not claim equivalence between the two.

\paragraph{Generalizability.}
The fingerprint characterization and replay study are specific to \Frontier's architecture,
cooling infrastructure, and workload mix.  Whether persistent thermal fingerprints of
comparable predictive accuracy and scheduling-capacity consequence exist on other GPU
supercomputers is an empirical question that cannot be answered from this study.

\paragraph{AI-assistance disclosure.}
This research used generative AI tools (Claude Code, Anthropic) for brainstorming,
methodological critique, code-generation assistance, analysis auditing, and
drafting/editing.  All study-design, preregistration, analysis, and interpretation
decisions were made and approved by the authors, who supervised or executed all
analyses, independently verified numerical claims against frozen artifacts, and
take full responsibility for all content.  AI tools participated in hypothesis and
comparison brainstorming, not only prose and code; the authors retained final
decision authority throughout.

%% ================================================================
\section{Conclusion}
\label{sec:conclusion}
%% ================================================================

We have shown that GPU nodes on the \Frontier supercomputer exhibit persistent thermal
fingerprints that are predictive across calendar years: fingerprints computed from 2024
telemetry prospectively predict 2025 node-level thermal positions with Pearson $r = 0.926$
and no-refit $R^2 = 0.848$.  A preregistered held-out constraint-aware replay study
demonstrates that fingerprint-aware placement yields a median conservative bracket
separation of $+12.48\pp$ over random placement (25/25 scenarios, Holm $p = 4.77\times10^{-7}$)
under a concurrent-load cabinet reference at the preregistered choice ratio.

Three independent controls---within-cabinet permutation, global permutation, and
structural leave-one-out position template---all show that individual node identity
contributes beyond cabinet membership and fixed rack geometry, within the scope of the
replay model.  Random spreading provides no statistically distinguishable benefit.
The direction of the FP advantage is preserved across three thermal threshold definitions;
the P99 magnitude is floor-limited and should not be interpreted as a larger effect.

Key open questions include: whether persistent fingerprints generalize to GPU
supercomputers with different cooling architectures; how fingerprint stability degrades
after hardware maintenance or node replacements; whether real-time deployment in a
production scheduler can realize the predicted capacity advantage; and how fingerprint
information should be weighted against competing placement objectives such as network
locality and power draw.

%% ================================================================
%% References
%% ================================================================

\bibliographystyle{ACM-Reference-Format}
\bibliography{thermalops}

%% ================================================================
%% Appendix
%% ================================================================

\appendix

\section{Preflight Conditions}

The following 16 conditions must all pass before scenario dispatch:
(1) Runner SHA256 matches frozen value;
(2) reference file SHA256 matches frozen value;
(3) denominator file SHA256 matches frozen value;
(4) git repository is tracked-clean with no uncommitted changes;
(5) \texttt{THERMALSHIFT\_DATA\_DIR} resolves to an existing parquet file;
(6--16) scenario-level checks including seed correctness, no pre-existing output files,
and simulation equivalence at $10^{-9}$ tolerance for one smoke-test scenario.

\section{Position-LOO Fallback Hierarchy}

\begin{description}
  \item[Level 1 (exact-position LOO):] All jobs at the same rack position excluding
    the evaluated job.  Applies to 1,990/2,000 nodes.
  \item[Level 2 (adjacent-level LOO):] Jobs at adjacent rack-level positions.
    No level-2 fallbacks were required.
  \item[Level 3 (same-cabinet LOO):] Jobs in the same cabinet.  Applies to 3 nodes.
  \item[Global fallback:] Fleet-wide average.  Applies to 7 nodes.
\end{description}

Seven global fallbacks were retained after an explicit pre-run ruling that their
position history was insufficient for a reliable level-1 or level-3 estimate.
The resulting template has Spearman $\rho = 0.9823$ correlation with the
unconstrained structural template.

\section{Binding Constraint Definitions}
\label{app:constraints}

\begin{description}
  \item[$B_\mathrm{ratio}$:] Thermal-incidence-ratio constraint.  The ratio of actual
    cabinet thermal incidence to $B_{\mathrm{ref},s}$ reaches the capacity boundary
    before the severity ratio does.  This is a thermal constraint (incidence dimension).

  \item[$S_\mathrm{ratio}$:] Thermal-severity-ratio constraint.  The ratio of actual
    severity to $S_{\mathrm{ref},s}$ reaches the boundary first.  Also a thermal
    constraint (severity dimension).
\end{description}

Both $B_\mathrm{ratio}$ and $S_\mathrm{ratio}$ are thermal constraints; the
distinction is which exposure dimension binds first.  At P90 and P95, FP crossings
are 100\% $B_\mathrm{ratio}$-driven.  At P99, 20/25 FP crossings shift to
$S_\mathrm{ratio}$, consistent with the hotter threshold definition placing more
weight on severity exposure before incidence saturates.

\section{Reproducibility}

All analysis scripts, frozen preregistration plans, result JSONs, and figure-generation
scripts are released at \url{https://github.com/K2alyan/thermalshift}
(archived: DOI 10.5281/zenodo.23020901).  Analysis plans are tied to specific commits
with SHA256-verified runner binaries.  The repository does not mirror raw \Frontier
telemetry; the underlying job metadata and per-node GPU temperature telemetry are
available through the OLCF Open Data Initiative under its access terms
(DOI: 10.13139/OLCF/3013979; acquisition instructions in \texttt{DATA.md}).

\end{document}
